% ===== §17 =====
\section{Room for the Core Without Possession}
\label{sec:ethicsgen}

\noindent
This section takes the second register of the limit's generativity, ethics, and establishes that the relation to the other founded on the limit is a giving that does not possess: a generating of the bond that leaves the other's unsymbolizable core its room, and that arrives thereby at the form older wisdom named, to generate and not possess. The section gathers the aesthetic measure of Part One and the making-room of Part Three into an ethical form, states that form and its ground, gives it as a claim, and marks its distance from an ethics of possession.

\subsection*{The other who cannot be exhausted}

The ethical relation begins where the other cannot be exhausted. This was established in Part One: the other faces me not as an instance of a kind but as this one, whose claim issues from what no predicate contains, and the fitting response to them is taken by a perception, an aesthetic measure, that no rule can replace. Were the other fully sayable, fully brought under description, there would be no ethical claim, only the management of a known quantity. The claim arises from what in the other exceeds every predicate, the singular core that no category reaches, and Part Three identified that core with the seat of the other's desire, the void around which their drive turns. Ethics is founded, then, not on the reaching of the other but on the acknowledgment of what cannot be reached, and its first act is to leave that core its room.

\subsection*{The aesthetic measure at work}

This is the aesthetic measure of Part One, now doing ethical work. The fitting response to another is one that senses where the sayable ends and does not press past it, that answers the other without presuming to have exhausted them, that keeps its response open to a singularity it perceives but cannot state. To make room for the other's core is not to withhold engagement but to engage in the mode that perception, not proposition, governs: to attend to the other closely enough to sense the edge of what may be said, and to hold there. The schooling of this perception, the cultivation the series has developed elsewhere, is the ethical discipline proper to a relation whose measure is aesthetic.

\subsection*{Giving without possession}

The form of this ethics is a giving that does not possess. To make room for the other's core is to be in relation to what one does not hold, to generate a bond with what one has not seized and cannot. The paper arrives here at a shape that an older tradition named, and names it as a conclusion reached on its own path rather than borrowed as an authority: to generate and not possess, to act and not hold onto, to foster and not rule. The ethical relation to the other's unsymbolizable core has exactly this form. It generates a bond and does not possess the other; it acts toward them and does not hold them fast; it fosters what the other is and does not rule over the core from which they desire. The arrival is not decorative. The account of Part Two forced it: since to possess the other's core, to close it, would be to extinguish the desire that makes the other a living source of the bond, an ethics that generated a lasting bond could only be an ethics that generated without possessing. The old formula names the only ethical form the structure permits.

\begin{claim}[Generation without possession]
The ethics grounded in the limit is a giving that does not possess: it generates a bond with the other while leaving the other's unsymbolizable core its room, declining to close what it cannot close without extinguishing the other's desire. To generate and not possess, to act and not hold onto, to foster and not rule, is not a borrowed maxim but the only ethical form the structure of the limit permits, since possession of the core would be the death of the desire that makes the other a living source of the bond.
\end{claim}

\subsection*{Against an ethics of possession}

This is the ethical face of the limit's generativity, and its distance from an ethics of possession is exact. An ethics that aimed at the full reaching of the other, the closing of every gap between them, would be an ethics of possession, and possession, here as throughout the argument, is extinction: the other possessed is the other's desire gone out. The ethics the limit grounds is the opposite of possession. It is the art of dwelling with what one cannot hold, of answering a singularity one cannot state, of generating a bond precisely by leaving open the core one will not close. What cannot be reached in the other is not the limit of ethics. It is the space in which an ethics worthy of the name becomes possible, and the space is kept open by the same refusal of closure that the poem practices in language and that justice, in the next section, practices in the common life.
